\documentclass[../../main.tex]{subfiles}
\begin{document}
\section{习题}\label{section:ch13-problems}

\begin{exercise}\label{exercise:ch13-01}
若 $(M,g)$ 是一个带边或不带边的 Riemann $n$ 流形，令 $UM\subseteq TM$ 为子集 $UM=\{(x,v)\in TM:|v|_g=1\}$，称为 $M$ 的\emph{单位切丛}。证明 $UM$ 是 $M$ 上的一个光滑纤维丛，其模型纤维为 $S^{n-1}$。
\end{exercise}

\begin{exercise}\label{exercise:ch13-02}
在命题\ref{proposition:ch13-existence}的证明中，我们用单位分解把局部定义的 Riemann 度量拼接起来得到一个整体度量。该证明的关键一步是验证这样得到的整体张量场正定。这一论证成功的关键在于：给定切空间上内积所成之集是所有对称 $2$-张量所成向量空间的一个凸子集。本题把这一构造推广到任意向量丛。设 $E$ 是带边或不带边光滑流形 $M$ 上的一个光滑向量丛，$V\subseteq E$ 是一个开子集，满足对每个 $p\in M$，$V$ 与纤维 $E_p$ 的交是凸且非空的。所谓“$V$ 的截面”，指 $E$ 的一个（局部或整体）截面，其像含于 $V$ 中。
\begin{enumerate}[(a)]
\item 证明存在 $V$ 的一个光滑整体截面。
\item 设 $\sigma:A\to V$ 是定义在闭子集 $A\subseteq M$ 上的一个 $V$ 的光滑截面。（这意味着 $\sigma$ 在 $A$ 的每一点的邻域内都可延拓为 $V$ 的一个光滑截面。）证明存在 $V$ 的一个光滑整体截面 $\tilde\sigma$，使其在 $A$ 上的限制等于 $\sigma$。证明：若 $V$ 包含 $E$ 的零截面的像，则 $\tilde\sigma$ 可被选成支集含于 $A$ 的任意预先给定的邻域内。
\end{enumerate}
（用于第~381、430 页。）
\end{exercise}

\begin{exercise}\label{exercise:ch13-03}
设 $M$ 是一个光滑流形。证明下列命题。
\begin{enumerate}[(a)]
\item 若 $M$ 上存在一个整体非零向量场，则存在一个整体光滑非零向量场。[提示：模仿定理\ref{theorem:ch06-21}的证明，把常数 $F(x^i)$ 换成坐标中的常系数向量场，并把绝对值换成某个 Riemann 度量下的范数。]
\item 若 $M$ 上存在一组线性无关的 $k$ 个向量场，则存在这样一组光滑向量场。
\end{enumerate}
\end{exercise}

\begin{exercise}\label{exercise:ch13-04}
令 $\bar g$ 表示 $S^n$ 上的圆度量。计算 $\bar g$ 在球极坐标下的坐标表示（见习题\ref{exercise:ch01-07}）。
\end{exercise}

\begin{exercise}\label{exercise:ch13-05}
设 $(M,g)$ 是一个 Riemann 流形。一个光滑曲线 $\gamma:J\to M$ 若满足 $|\gamma'(t)|_g\equiv1$，就称为一条\emph{单位速率曲线}。证明：每条速度处处非零的光滑曲线都有一个单位速率重新参数化。（用于第~335 页。）
\end{exercise}

\begin{exercise}\label{exercise:ch13-06}
证明每个 Riemann $1$ 流形都是平坦的。[提示：用习题\ref{exercise:ch13-05}。注意这蕴含 $S^1$ 上的圆度量是平坦的！]
\end{exercise}

\begin{exercise}\label{exercise:ch13-07b}
证明平坦度量的乘积仍是平坦的。
\end{exercise}

\begin{exercise}\label{exercise:ch13-08}
令 $\mathbb{T}^n=S^1\times\cdots\times S^1\subseteq\mathbb{C}^n$，令 $g$ 为 $\mathbb{T}^n$ 上由 $\mathbb{C}^n$（等同于 $\mathbb{R}^{2n}$）上 Euclid 度量诱导的度量。证明 $g$ 是平坦的。
\end{exercise}

\begin{exercise}\label{exercise:ch13-09}
令 $H\subseteq\mathbb{R}^3$ 为螺旋面（即例\ref{example:ch13-helicoid}中嵌入 $F:\mathbb{R}^2\to\mathbb{R}^3$ 的像），令 $C\subseteq\mathbb{R}^3$ 为悬链面，即由曲线 $\gamma(t)=(\cosh t,t)$ 生成的旋转曲面。证明 $H$ 局部等距于 $C$，但不整体等距。
\end{exercise}

\begin{exercise}\label{exercise:ch13-10}
证明 Euclid 空间中两点之间的最短路是一条直线段。更确切地说，对 $x,y\in\mathbb{R}^n$，令 $\gamma:[0,1]\to\mathbb{R}^n$ 为曲线段 $\gamma(t)=x+t(y-x)$，证明：任何其他从 $x$ 到 $y$ 的分段光滑曲线段 $\widetilde\gamma$，除非它是 $\gamma$ 的一个重新参数化，都满足 $L_{\bar g}(\widetilde\gamma)>L_{\bar g}(\gamma)$。[提示：首先考虑 $x,y$ 都在 $x^1$-轴上的情形。]（用于第~338 页。）
\end{exercise}

\begin{exercise}\label{exercise:ch13-11}
令 $M=\mathbb{R}^2\setminus\{0\}$，令 $g$ 为 Euclid 度量 $\bar g$ 在 $M$ 上的限制。证明存在点 $p,q\in M$，使得 $M$ 中不存在从 $p$ 到 $q$、满足 $L_g(\gamma)=d_g(p,q)$ 的分段光滑曲线段 $\gamma$。
\end{exercise}

\begin{exercise}\label{exercise:ch13-12}
把 $\mathbb{R}^n$ 视为带 Euclid 度量 $\bar g$ 的 Riemann 流形。
\begin{enumerate}[(a)]
\item 设 $U,V\subseteq\mathbb{R}^n$ 是连通开集，$\varphi,\psi:U\to V$ 是 Riemann 等距，且对某个 $p\in U$ 满足 $\varphi(p)=\psi(p)$、$\mathrm{d}\varphi_p=\mathrm{d}\psi_p$。证明 $\varphi=\psi$。[提示：先用习题\ref{exercise:ch13-10}证明 $\varphi,\psi$ 把直线变成直线。]
\item 证明：例\ref{example:ch07-euclidean-group}中所述 $E(n)$ 在 $\mathbb{R}^n$ 上的作用所给出的、从 $\mathbb{R}^n$ 到自身的映射集合，就是 $(\mathbb{R}^n,\bar g)$ 的整个 Riemann 等距群。
\end{enumerate}
\end{exercise}

\begin{exercise}\label{exercise:ch13-13}
设 $(M,g)$ 是一个 Riemann 流形。$M$ 上的一个光滑向量场 $V$ 若其流按 $g$ 的等距作用，就称为 $g$ 的一个\emph{Killing 向量场}（以十九世纪末二十世纪初德国数学家 Wilhelm Killing 命名）。
\begin{enumerate}[(a)]
\item 证明 $M$ 上所有 Killing 向量场构成 $\mathfrak X(M)$ 的一个 Lie 子代数。[提示：见推论\ref{corollary:ch09-flow-derivatives}(c)。]
\item 证明 $M$ 上一个光滑向量场 $V$ 是 Killing 向量场当且仅当它在每个光滑局部坐标图中满足下列方程：
\begin{align*}
V^k\frac{\partial g_{ij}}{\partial x^k}+g_{jk}\frac{\partial V^k}{\partial x^i}+g_{ik}\frac{\partial V^k}{\partial x^j}=0.
\end{align*}
\end{enumerate}
\end{exercise}

\begin{exercise}\label{exercise:ch13-14}
令 $\mathcal K\subseteq\mathfrak X(\mathbb{R}^n)$ 表示关于 Euclid 度量的 Killing 向量场所成 Lie 代数（见习题\ref{exercise:ch13-13}），令 $\mathcal K_0\subseteq\mathcal K$ 表示由在原点处消失的向量场所成的子空间。
\begin{enumerate}[(a)]
\item 证明映射
\begin{align*}
V\mapsto\left(\frac{\partial V^i}{\partial x^j}(0)\right)
\end{align*}
是从 $\mathcal K_0$ 到 $\mathfrak o(n)$ 的一个单射线性映射。[提示：若 $V$ 在此映射的核中，$\theta$ 是它的流，证明线性映射 $\mathrm{d}(\theta_t)_0:T_0\mathbb{R}^n\to T_0\mathbb{R}^n$ 与 $t$ 无关，并利用习题\ref{exercise:ch13-12}(a)的结果。]
\item 证明下列向量场构成 $\mathcal K$ 的一组基：
\begin{align*}
\frac{\partial}{\partial x^i},\quad 1\leqslant i\leqslant n;\qquad
x^i\frac{\partial}{\partial x^j}-x^j\frac{\partial}{\partial x^i},\quad 1\leqslant i<j\leqslant n.
\end{align*}
\end{enumerate}
\end{exercise}

\begin{exercise}\label{exercise:ch13-15}
设 $(M,g)$ 是一个 Riemann 流形，令 $\hat g$ 为 $M$ 与 $\mathbb{R}$ 上 Euclid 度量在 $M\times\mathbb{R}$ 上确定的乘积度量。令 $X=0\oplus\mathrm{d}/\mathrm{d}t$ 为 $M\times\mathbb{R}$ 上由 $M$ 上零向量场与 $\mathbb{R}$ 上标准坐标向量场 $\mathrm{d}/\mathrm{d}t$ 确定的乘积向量场（见习题\ref{exercise:ch08-17}）。证明 $X$ 是 $(M\times\mathbb{R},\hat g)$ 的一个 Killing 向量场。
\end{exercise}

\begin{exercise}\label{exercise:ch13-16}
设 $g=f(t)\,\mathrm{d}t^2$ 是 $\mathbb{R}$ 上的一个 Riemann 度量。证明 $g$ 完备当且仅当下列两个反常积分都发散：
\begin{align*}
\int_0^\infty \sqrt{f(t)}\,\mathrm{d}t,\qquad
\int_{-\infty}^0 \sqrt{f(t)}\,\mathrm{d}t.
\end{align*}
\end{exercise}

\begin{exercise}\label{exercise:ch13-17}
设 $M$ 是一个连通的非紧光滑流形，$g$ 是 $M$ 上的一个 Riemann 度量。证明存在一个正函数 $h\in C^\infty(M)$，使 Riemann 度量 $\widetilde g=hg$ 完备。利用这一点证明每个连通光滑流形都容许一个完备 Riemann 度量。[提示：令 $f:M\to\mathbb{R}$ 为一个穷竭函数，并证明 $h$ 可选成使 $f$ 在 $\widetilde g$-有界集上有界。]
\end{exercise}

\begin{exercise}\label{exercise:ch13-18}
设 $(M,g)$ 是一个连通 Riemann 流形，$S\subseteq M$ 是一个连通嵌入子流形，$\widetilde g$ 是 $S$ 上的诱导 Riemann 度量。
\begin{enumerate}[(a)]
\item 证明对 $p,q\in S$，有 $d_{\widetilde g}(p,q)\geqslant d_g(p,q)$。
\item 证明：若 $(M,g)$ 完备且 $S$ 正常嵌入，则 $(S,\widetilde g)$ 完备。
\item 利用 (b) 以及 Whitney 嵌入定理证明（不引用命题\ref{proposition:ch13-existence}或习题\ref{exercise:ch13-17}）：每个连通光滑流形都容许一个完备 Riemann 度量。
\end{enumerate}
\end{exercise}

\begin{exercise}\label{exercise:ch13-19}
下例表明习题\ref{exercise:ch13-18}(b) 的逆不成立。由 $F(t)=((\mathrm{e}^t+1)\cos t,(\mathrm{e}^t+1)\sin t)$ 定义 $F:\mathbb{R}\to\mathbb{R}^2$。证明 $F$ 是一个非正规嵌入，但 $\mathbb{R}$ 关于从 $\mathbb{R}^2$ 上 Euclid 度量诱导的度量是完备的。
\end{exercise}

\begin{exercise}\label{exercise:ch13-20}
考虑由下式定义的嵌入 $F_1,F_2,F_3,F_4:\mathbb{R}^2\to\mathbb{R}^3$：
\begin{align*}
F_1(u,v)&=(u,v,0);\\
F_2(u,v)&=(u,\mathrm{e}^v,0);\\
F_3(u,v)&=(u,v,u^2+v^2);\\
F_4(u,v)&=\left(\frac{2u}{u^2+v^2+1},\frac{2v}{u^2+v^2+1},\frac{u^2+v^2-1}{u^2+v^2+1}\right).
\end{align*}
对每个 $i$，令 $g_i=F_i^*\bar g$。对 Riemann 流形 $(\mathbb{R}^2,g_1)$ 直至 $(\mathbb{R}^2,g_4)$ 中的每一个，回答下列问题：它有界吗？它完备吗？它平坦吗？证明你的答案正确。[提示：若分析 $g_4$ 有困难，看习题\ref{exercise:ch01-07}。]
\end{exercise}

\begin{exercise}\label{exercise:ch13-21}
设 $(M,g)$ 是一个 Riemann 流形，$f\in C^\infty(M)$，$p\in M$ 是 $f$ 的一个正则点。
\begin{enumerate}[(a)]
\item 证明：在所有单位向量 $v\in T_pM$ 中，方向导数 $vf$ 在 $v$ 与 $\operatorname{grad}f|_p$ 同向时最大，且 $\operatorname{grad}f|_p$ 的长度等于该方向上方向导数的值。
\item 证明 $\operatorname{grad}f|_p$ 正交于 $f$ 过 $p$ 的水平集。
\end{enumerate}
（用于第~391 页。）
\end{exercise}

\begin{exercise}\label{exercise:ch13-22}
对任意带边或不带边光滑流形 $M$，证明向量丛 $TM$ 与 $T^*M$ 在 $M$ 上光滑同构。[注：习题\ref{exercise:ch11-18}表明这个同构不能是自然的，即不存在这样一条规则：给每个光滑流形 $M$ 指派一个丛同构 $\Phi_M:TM\to T^*M$，使得对每个微分同胚 $F:M\to N$，两个丛同构 $\Phi_M$ 与 $\Phi_N$ 满足 $\Phi_M=\mathrm{d}F^*\circ\Phi_N\circ\mathrm{d}F$。]
\end{exercise}

\begin{exercise}\label{exercise:ch13-23}
$S^2$ 上是否存在一个恰好在一点处消失的光滑余向量场？
\end{exercise}

\begin{exercise}\label{exercise:ch13-24}
设 $M$ 是一个紧的光滑 $n$ 流形，$f$ 是 $M$ 上一个光滑实值函数，只有有限多个临界点 $\{p_1,\dots,p_k\}$，相应的临界值为 $\{c_1,\dots,c_k\}$，标号使 $c_1\leqslant\cdots\leqslant c_k$。对任意 $a<b\in\mathbb{R}$，定义 $M_a=f^{-1}(a)$、$M_{[a,b]}=f^{-1}([a,b])$、$M_{(a,b)}=f^{-1}((a,b))$。若 $a,b$ 是正则值，注意 $M_a,M_b$ 是 $M$ 中的嵌入超曲面；$M_{(a,b)}$ 是一个开子流形；而 $M_{[a,b]}$ 由命题\ref{proposition:ch05-regular-domain}是一个正则区域（图\ref{figure:ch13-morse}）。
\begin{enumerate}[(a)]
\item 在 $M$ 上取一个 Riemann 度量 $g$，令 $X$ 为向量场 $X=\operatorname{grad}f/|\operatorname{grad}f|_g^2$（在 $M\setminus\{p_1,\dots,p_k\}$ 上），令 $\theta$ 为 $X$ 的流。证明只要 $\theta_t(p)$ 有定义，就有 $f(\theta_t(p))=f(p)+t$。
\end{enumerate}

\begin{figure}[htbp]
\centering
\includegraphics[width=0.6\textwidth]{fig13-4-morse.png}
\caption{习题\ref{exercise:ch13-24}的设定}
\label{figure:ch13-morse}
\end{figure}

\begin{enumerate}[(b)]
\item 令 $[a,b]\subseteq\mathbb{R}$ 为一个不含 $f$ 的临界值的紧区间。证明 $\theta$ 限制为从 $[0,b-a]\times M_a$ 到 $M_{[a,b]}$ 的一个微分同胚。
\end{enumerate}
[注：这一结果表明 $M$ 可分解为一些较简单的“积木块”之并——乘积集 $M_{[c_i+\varepsilon,c_{i+1}-\varepsilon]}\approx I\times M_{c_i+\varepsilon}$，以及临界点附近的邻域 $M_{(c_i-\varepsilon,c_i+\varepsilon)}$。这就是 Morse 理论的出发点，它是微分几何对拓扑最深刻的应用之一。下一步将分析 $f$ 在每个临界点附近的行为，并利用这一分析确定当穿过临界水平时水平集在拓扑上究竟如何变化。精彩的引论见\cite{Mil63}。]
\end{exercise}

\begin{exercise}\label{exercise:ch13-25}
设 $M$ 是一个光滑流形，它容许一个没有临界点的正常光滑函数 $f:M\to\mathbb{R}$。证明 $M$ 微分同胚于 $N\times\mathbb{R}$，其中 $N$ 是某个紧光滑流形。[提示：令 $X=\operatorname{grad}f/|\operatorname{grad}f|_g$，它关于 $M$ 上某个 Riemann 度量定义。证明 $X$ 是完备的，并利用它的流形扩张来定义这个微分同胚。]
\end{exercise}
\end{document}
